\section{Scalability and many capacities}
\label{sec:scalability}

\generated{table_scalability}

\begin{figure}[tbp]
  \centering
  \IfFileExists{generated/scaling_scalability_parametric.dat}{%
    \begin{tikzpicture}
      \begin{axis}[xlabel={vertices $n$}, ylabel={median seconds}, xmode=log, ymode=log,
                   legend pos=north west, width=0.62\textwidth, height=0.42\textwidth]
        \IfFileExists{generated/scaling_scalability_forest-tuned.dat}{%
          \addplot table[x=nodes, y=seconds] {generated/scaling_scalability_forest-tuned.dat};
          \addlegendentry{\texttt{forest-tuned}}}{}
        \IfFileExists{generated/scaling_scalability_dinkelbach.dat}{%
          \addplot table[x=nodes, y=seconds] {generated/scaling_scalability_dinkelbach.dat};
          \addlegendentry{\texttt{dinkelbach}}}{}
        \addplot table[x=nodes, y=seconds] {generated/scaling_scalability_parametric.dat};
        \addlegendentry{\texttt{parametric}}
      \end{axis}
    \end{tikzpicture}}{\small\itshape Scalability data not available: the
    scalability campaign has not been analysed in this build of the report.}
  \caption{Median wall time against the number of vertices on the generated
  families whose structure is held fixed while the size varies. Only runs that
  finished and were verified appear; a configuration whose curve stops has hit
  the time limit beyond that size.}
  \label{fig:scaling}
\end{figure}

\subsection{One capacity against many}

\generated{table_capacities}

The canonical positive path is built once and answers any capacity: locating the
critical block costs $O(\log k)$ in the number of blocks, but exporting a
dense solution vector costs at least $O(n)$ per query. The campaign therefore
separates the cost of the path from the cost of the queries, and reports the
number of queries at which building the whole path pays for itself against
solving each capacity independently.
